\subsection{Retained reconstruction with a fixed homothety origin}
In the stationary experiment, a forward attempt at nominal center $x$
starts at $x+Z$, with an independent compass displacement, and a reverse
attempt starts at $x+Z+a$ with displacement $-a$. The draw $Z$ has an
unknown density on an unknown convex body $K$. The second setting has
footprint $rK$, where the unknown ratio satisfies known bounds
$1+g_*\le r\le R_*$. We take the first operating footprint as the
reference unit of dilation. Both supports have fixed positive size;
they do not contract as accuracy increases. Homothety is about the
specified laboratory origin. The numerical inball, envelope, curvature
and $C^{6,\beta}$ bounds and a lower boundary-mass condition are given.
In particular, for fixed $b_0>0$ and $\gamma\ge0$ the first density obeys
$j(z)\ge b_0\dist(z,\partial K)^\gamma$ almost everywhere in $K$.
Section~\ref{sec:unknown-footprint} gives the detailed geometric priors.

Choose once a dyadic compass step so that the largest footprint
diameter plus $2t$ is less than $d_0$. This stricter design leaves a
positive collar for the translated commands used below. It is a
prior-dependent choice of the same four-direction pooled apparatus.
No direction-resolved response, launch position, occupancy bit or
obstacle area is observed.

The retained fixed-origin constructive conclusion is
Theorem~\ref{thm:unk-scale-finite}. It recovers $K$, $r$ and the original
table to laboratory $C^2$ error $C\nu$, and, on $\mathcal B_\eta$, the
primitive discrete data with probability at least $1-\delta$. Put
\begin{equation}\label{eq:new-upper-exponent}
 A_s(\gamma)=\frac{(\gamma+3/2)s+1}{s-2},\qquad
 L_\nu=\log(C/\nu),\qquad L_{\nu,\delta}=\log(C/(\nu\delta)).
\end{equation}
The sufficient attempted-bit and command-description bounds are
\begin{align}
 N_\nu&\le C\nu^{-A_s(\gamma)}L_\nu L_{\nu,\delta},
                                              \label{eq:main-new-attempts}\\
 \mathcal D_\nu&\le C\bigl\{\nu^{-1/(s-2)}L_\nu
                      +\nu^{-2}L_{\nu,\delta}\bigr\}L_{\nu,\delta}.
                                              \label{eq:main-new-control}
\end{align}
Here $\mathcal D_\nu$ is the binary length of the commanded centers,
setting indices and repetition counts. The finest dyadic nominal mesh
is $O(\nu^{s/(s-2)})$. These statements price the input descriptions
alongside the observations; physical manufacture, travel, homothety
certification and arithmetic running time remain separate resources.
All attempted preparations, including solid starts and misses, are
included in $N_\nu$. The numerical priors are known with certified
enclosures; neither $r$ nor a support-function evaluation oracle for
$K$ is supplied.

There are two new steps. First, near each recovered coarse boundary,
a finite conjunction of shifted pooled collision tests has a zero
probability outside and a uniformly positive probability at depth $e$.
Proposition~\ref{prop:rare-query} proves that its cost is
$O(e^{-(\gamma+3/2)}\log(1/\varepsilon))$. This avoids the squared
inverse-signal cost of the signed occupation query. Second, the integral of a pooled forward mean measures the sum of two
obstacle widths. A linear identity with the two positivity supports
then determines the unknown homothety ratio. Its $O(\nu^{-2})$
sampling cost is absorbed
by \eqref{eq:main-new-attempts}, since $6<s\le7$ and
$A_s(\gamma)>2$. As a second exact argument, a finite killed-walk
adjoint recovers one occupation integral, equal to the obstacle area;
a mixed-area identity determines the ratio from the two reciprocal
differences alone, under a weaker separation condition.

For $\gamma=0$, the common-noise converse in
Theorem~\ref{thm:stationary-lower} gives lower power $(s+1)/(s-2)$,
whereas \eqref{eq:main-new-attempts} has power $(3s/2+1)/(s-2)$.
Thus the difference of the powers is $s/(2(s-2))$; it was $2s/(s-2)$
for the retained signed-query construction. This earlier comparison alone does not identify the minimax exponent.
The sharp common-disk result is now Corollary~\ref{cor:stationary-minimax}. The converse permits arbitrary-precision
adaptive short commands under a single known uniform-disk law and
charges worst-case expected attempts. The finite period assertion
continues to use the known positive nonperiod-patch margin. Exact
geometric recovery and pointwise eventual period recognition are
stated separately from that uniform finite assertion.

\subsection{Location of the retained stationary arguments}
Section~\ref{sec:stationary} develops the stationary occupation inverse
and its quantitative support geometry. Section~\ref{sec:unknown-footprint}
separates two footprints when their scales are known.
Section~\ref{sec:rare-stationary} proves the faster boundary query;
Section~\ref{sec:unknown-scale} identifies an unknown ratio by integrated
collision width and, independently, by occupation area.
Section~\ref{sec:stationary-information} proves the common-noise
converse. The appendices contain the retained localized statements,
coarse reconstruction, smooth interpolation, period tests, physical
packing, finite controls and calibration converse, with their proofs.
